\def\stat{zat-kis}



\def\tit{НЕКОТОРЫЕ ПОДХОДЫ К ФОРМИРОВАНИЮ НОРМАТИВНО-ТЕХНИЧЕСКОЙ 

БАЗЫ В~ЧАСТИ ТРЕБОВАНИЙ К~АРХИТЕКТУРНОМУ ПОСТРОЕНИЮ 

ИНФОРМАЦИОННЫХ СИСТЕМ ОРГАНИЗАЦИЙ~--- УЧАСТНИКОВ ЕДИНОГО 

ИНФОРМАЦИОННОГО ПРОСТРАНСТВА РОССИИ}



\def\titkol{Некоторые подходы к~формированию нормативно-технической 

базы} % в~части требований к~архитектурному построению  информационных систем организаций~--- участников единого  информационного пространства России}



\def\autkol{А.\,А. Зацаринный, Э.\,В.~Киселев}





\def\aut{А.\,А. Зацаринный$^1$, Э.\,В.~Киселев$^2$}



\titel{\tit}{\aut}{\autkol}{\titkol}



%{\renewcommand{\thefootnote}{\fnsymbol{footnote}}

%\footnotetext[1]

%{Работа выполнена при финансовой поддержке  РФФИ (проект 15-07-02244).}}



\renewcommand{\thefootnote}{\arabic{footnote}}

\footnotetext[1]{Федеральный исследовательский

центр <<Информатика и~управление>> Российской академии наук, azatsarinny@ipiran.ru}

\footnotetext[2]{Институт проблем информатики Федерального исследовательского

центра <<Информатика и~управление>> Российской академии наук, ekiselev @ipiran.ru}



\label{st\stat}



                  \thispagestyle{headings}

 



\Abst{Статья продолжает серию публикаций, посвященных формированию  

нор\-ма\-тив\-но-тех\-ни\-че\-ской базы для создания информационных систем (ИС)

организаций~--- 

участников единого информационного пространства Российской Федерации (ЕИП РФ). 

Представлены обобщенные требования к~системам, включая: базовые принципы 

разработки проектных решений по\-стро\-ения и~создания систем; необходимые свойства 

систем; состав необходимого документационного обеспечения систем; направления 

возможного развития ИС; требования по использованию 

математического и~натурного моделирования при создании и~эксплуатации систем; 

требования по автоматизированному управлению функционированием систем; 

соблюдение принципа обеспечения информационной безопасности (ИБ) при по\-стро\-ении 

систем; выполнение принципа экономической целесообразности и~обоснованности 

финансовых затрат на развертывание и~поддержание работоспособности и~актуализации 

систем; обязательность соблюдения принципов преемственности, направленной эволюции 

и максимального использования вложенных ин\-вес\-ти\-ций в~системы; акцентирование 

внимания на особую значимость решения организационных вопросов при 

проектировании, создании и~эксплуатации систем в~течение их жизненного цикла. 

Положения данной статьи ориентированы на использование в~качестве  

научно-методической основы при формировании нормативно-технической базы 

построения ИС многочисленных организаций~--- участников ЕИП 

РФ.}



\KW{единое информационное пространство Российской Федерации (ЕИП РФ); 

организация~--- участник ЕИП РФ (УЕИП); системный подход; архитектура единой 

информационной системы УЕИП; защищенные информационные ресурсы; 

централизация; интеграция; информационная без\-опас\-ность; телекоммуникационное 

обеспечение}



\DOI{10.14357\!/\!08696527150311}



                     %\vspace*{-12pt}







\vskip 12pt plus 9pt minus 6pt





      \thispagestyle{headings} 

    

\section{Вводные положения}



     Ранее отмечалось~[1, 2], что потенциально участниками единого 

информационного пространства Российской Федерации (далее~--- УЕИП) 

могут стать многочисленные и~разнообразные организации. 

Информатизацию УЕИП следует предусматривать на основе системного 

подхода применительно к~архитектурам единых ИС (ЕИС)

УЕИП. Такой подход предусматривает \textit{создание стратегии УЕИП 

в~сфере своей деятельности на заданный период времени} (далее~--- 

стратегия). Стратегия~--- определенная программа действий, разработанная 

руководством для поддержания эффективного функционирования и~развития 

УЕИП. При разработке стратегии основное внимание следует уделять 

формированию планов, которые необходимы для контроля эффективности 

достижения стратегических ориентиров в~сфере деятельности УЕИП на 

заданный период времени. По существу, стратегия~--- это генеральная 

программа развития организации, определяющая стратегические задачи, их 

приоритеты, методы привлечения и~распределения ресурсов и~последовательность 

шагов по достижению стратегических целей 

и~в~наибольшей степени соответствующая сложившемуся состоянию внутренней 

и внешней среды для УЕИП.

     

     Стратегия должна включать перечень взаимоувязанных мероприятий 

по достижению заданных целевых показателей деятельности УЕИП, иметь 

упреждающий характер за счет механизмов адаптации в~условиях изменений 

в~сфере своей деятельности.

     

     Стратегия должна:

     \begin{itemize}

\item содержать конкретные цели, достижение которых является решающим 

для достижения заданных показателей в~соответствующей сфере 

деятельности;

     \item поддерживать рациональную инициативу на соответствующем 

уровне (федеральном, региональном, местном);

     \item концентрировать главные усилия в~нужное время в~нужном месте;

\item  предусматривать гибкость поведения и~принятия решений для 

достижения требуемого результата при минимуме ресурсов;

     \item обозначать скоординированное руководство на заданных уровнях 

(федеральном, региональном, муниципальном);

     \item предполагать корректное расписание действий на заданных 

уровнях (федеральном, региональном, муниципальном);

     \item обеспечивать гарантированные ресурсы для работы в~сфере своей 

деятель\-ности.

     \end{itemize}

     

     Стратегия должна выступать в~качестве нормативного инструмента 

обоснования, выработки и~реализации долгосрочных целей в~сфере своей 

деятельности, а также как средство связи УЕИП с~<<внешней средой>>: 

с~другими организациями~--- участниками ЕИП РФ различных уровней, 

с~физическими лицами~--- гражданами РФ, при необходимости 

с~соответствующими органами зарубежных государств (согласно 

установленному нормативно-правовому порядку). 

     

     \textit{Стратегия должна включать две части}:

     \begin{enumerate}[(1)]

     \item  собственно стратегию деятельности и~развития УЕИП 

(применительно к~организационной сфере своей деятельности);

     \item стратегию информатизации организации~--- УЕИП на основе 

положений первой части.

     \end{enumerate}

     

\section{Основные требования к~архитектурному построению единых 

информационных систем организаций~--- участников единого 

информационного пространства России}

     

     1.\  В соответствии с~положениями выработанной стратегии для 

организации~--- участника ЕИП РФ при информатизации необходимо 

\textit{решать одновременно двуединую укрупненную задачу}:

     \begin{enumerate}[(1)]

     \item ликвидировать внутри пространства данной 

конкретной организации пресловутую <<лоскутную>> информатизацию 

путем консолидации и~эволюционного создания ЕИС~[1, 2];

     \item посредством этой же ЕИС подключить и~органично ввести данную организацию в~ЕИП РФ на 

правах полноценного участника.

     \end{enumerate}

     

     Решение такой двуединой задачи позволит обеспечить эффективную 

информатизацию высокого, качественно нового уровня для каждой 

конкретной организации. Действительно, 

ЕИС~--- это инфраструктура организации, которая поддерживает 

решение, прежде всего, совокупности собственных актуальных задач. При 

этом такая инфраструктура обеспечивает достижение своих целей, а~именно:

     \begin{itemize}

     \item  ЕИС охватывает все физические 

объекты, объединяет на качественно новом уровне функционал всей 

совокупности существующих информационных систем. Таким образом 

создается единое внутреннее пространство своей организации. Единая 

ИС становится основой, единственной 

системообразующей информационной инфраструктурой, на базе и~в рамках 

которой конструируются все другие подсистемы/компоненты и~посредством 

которой решается совокупный функционал для своей организации. При этом 

ставится задача минимизации совокупной стоимости владения единой 

информационной системой;

     \item одновременно ЕИС берет на себя 

согласование, подключение и~органичное введение данной организации в~ЕИП РФ как полноценного участника.

     \end{itemize}

     

     \smallskip

     

     2.\ Информатизация многочисленных УЕИП, построенная на основе 

системного подхода к~формированию нормативно-технической базы 

и~выработанных соответствующих стратегий, напрямую затрагивает 

архитектурное построение таких ИС~[3--7], в~том 

чис\-ле предусматривается \textit{обобщение и~соответствующее 

типизирование архитектур ИС}. 

     

     В свою очередь, типизация архитектур предварительно предполагает 

обобщение определенных групп функций организаций~--- участников ЕИП 

РФ. Для примера можно рассмотреть определенное обобщение функций 

УЕИП в~сферах государственного управления России. Такие 

\textit{обобщенные функции} включают:

     \begin{itemize}

     \item выработку и~реализацию государственной политики;

     \item нормативно-правовое регулирование;

     \item правоприменительные функции;

     \item государственный контроль (надзор)~--- федеральный или 

региональный;

     \item предоставление (исполнение) государственных услуг.

     \end{itemize}

     

     Применительно к~эффективной информатизации УЕИП 

\textit{государственное управление в~сфере деятельности для конкретного 

участника} в~обобщенном виде охватывает следующие процессы: 

     \begin{itemize}

     \item  \textit{сбор и~обработку (анализ) информации}, относящейся 

     к~сфере деятельности организации; 

     \item  \textit{прогнозирование}, т.\,е.\ научное предвидение изменений 

     в~развитии ка\-ких-ли\-бо явлений или процессов в~сфере деятельности на основе 

объективных данных и~достижений науки; 

     \item \textit{планирование}, т.\,е.\ определение направлений, целей 

управленческой деятельности и~соответствующих способов, средств 

достижения этих целей; 

      \item \textit{собственно организацию}, т.\,е.\ формирование системы 

управления, упорядочение управленческих отношений между субъектами 

и~объектами управления, определение прав и~обязанностей, структуры органов, 

подчиненных организаций, подбор и~расстановка кадров и~т.\,д.; 

     \item  \textit{регулирование или распорядительство}, т.\,е.\ 

установление режима в~сфере деятельности по достижению целей и~задач 

управления, регулирование поведения управляемых объектов, подготовка 

директив, указаний, предписаний и~др.; 

     \item \textit{координацию и~взаимодействие}, осуществляемые для 

достижения общих целей управления в~сфере деятельности; 

     \item \textit{контроль и~учет}, состоящие в~определении степени 

соответствия фактического состояния объекта управления заданному.

     \end{itemize}

     

     В целом государственное управление в~сфере деятельности 

многочисленных УЕИП может охватывать совокупность сфер 

государственного управления федерального, регионального и~

муниципального уровней.

     

     Естественно, что в~соответствии со сферами их деятельности 

     у~различных УЕИП самые разнообразные информационные ресурсы (ИР), 

которыми они обмениваются между собой как внутри своих организаций, так 

в рамках ЕИП РФ. Под ИР сфер государственного управления обобщенно 

понимается совокупность данных, представляющих ценность для этих сфер 

и~выступающих в~качестве материальных ресурсов. К~ним относятся файлы 

данных, документы, тексты, таблицы, графики, аудиоинформация, 

видеоинформация, специфическая информация (например, 

дактилоскопическая идентификационная) и~др. Соответственно, необходимы 

ЕИС УЕИП как структуры, обеспечивающие 

хранение, поиск и~выдачу информации по запросам пользователей различных 

сфер деятельности, включая население РФ и~иностранных граждан.

     

     Правильное и~четкое выполнение различными УЕИП возложенных на 

них функций и~решение сопутствующих задач невозможно без эффективной 

информатизации собственной сферы деятельности на высоком, качественно 

новом уровне применительно к~каждой организации~--- участнику 

ЕИП~\cite{7-zk}. Данная ситуация требует рассмотрения многочисленных 

вопросов исключительно на базе системного подхода.

     

     В~[1, 2] определено, что организации-участники могут входить в~ЕИП 

РФ посредством систем двух видов: ЕИС 

и~единая ин\-фор\-ма\-ци\-он\-но-те\-ле\-ком\-му\-ни\-ка\-ци\-он\-ная сис\-те\-ма 

(ЕИТКС). Отличия 

только в~том, что в~ЕИТКС по сравнению с~ЕИС более существенную роль 

играет телекоммуникационная составляющая. 

     

     Выбор архитектуры и~создание ЕИС (ЕИТКС) конкретной УЕИП на 

основе требований системного подхода и~стратегии должны обеспечить 

\textit{эволюционный переход от существующей <<лоскутной>> 

информатизации с~постепенной заменой используемых ИС} в~этой 

организации.

     

     \smallskip

     

     3. Архитектура ЕИС (ЕИТКС) каждой УЕИП в~общем случае должна 

быть нацелена на обеспечение информационной, аналитической, 

документационной, инструментальной и~технологической поддержки 

принятия решений и~выполнения основных функций всего спектра сферы 

деятельности для конкретной организации применительно 

к~соответствующему уровню~--- федеральному, региональному, 

территориальному (муниципальному). 

     

     Ниже излагаются \textit{требования к~архитектурному построению 

обобщенной, типизированной архитектуры ЕИС (ЕИТКС)  

организации~--- участника ЕИП РФ}.

     

     \smallskip

     

     4. В состав архитектуры ЕИС (ЕИТКС) организации~--- участника ЕИП 

РФ в~общем случае могут быть включены \textit{три составные части~--- 

три сегмента}:

     \begin{enumerate}[(1)]

     \item \textit{открытый сегмент (О-сегмент)}, который оперирует 

открытой информацией и~аккумулирует открытые ИР. 

В рамках данного сегмента создается открытый контур ЕИС (ЕИТКС) 

     (О-кон\-тур~\cite{2-zk}); 

     \item \textit{конфиденциальный сегмент (К-сегмент)}, оперирующий 

конфиденциальной информацией и~аккумулирующий конфиденциальные 

ИР. В рамках данного сегмента создается 

конфиденциальный контур ЕИС (ЕИТКС)\linebreak 

(К-кон\-тур~\cite{2-zk}); 

     \item \textit{закрытый сегмент (З-сегмент)}, оперирующий только 

закрытой информацией. В~рамках данного сегмента создается закрытый 

контур ЕИС (ЕИТКС) (З-контур~\cite{2-zk}).

     \end{enumerate}

     

     Таким образом, \textit{ЕИС (ЕИТКС) в~общем случае фактически 

будет состоять из трех автоматизированных ИС

(АИС), которые представляют три перечисленных сегмента}. Между АИС 

сегментами посредством шлюзов создается слабая связность внутри ЕИС 

(ЕИТКС) конкретной УЕИП, допускающая при необходимости 

однонаправленную передачу~\cite{2-zk}: только из открытого контура 

в~конфиденциальный контур, только из конфиденциального контура 

в~закрытый контур. Открытый сегмент через собственный шлюз получает от 

доступных источников открытую первичную информацию.

     

     Число необходимых сегментов ЕИС (ЕИТКС) зависит от потребностей 

конкретной УЕИП.

     

     \smallskip

     

     5. Предлагается выделить следующие \textit{обобщенные требования 

системного подхода} к~построению, созданию и~эксплуатации ЕИС (ЕИТКС) 

организации~--- участника ЕИП РФ:

     \begin{itemize}

     \item[(а)] \textit{общность использования}: реализация ЕИС (ЕИТКС) 

должна обеспечивать взаимодействие неопределенного круга ИС~--- как 

собственных, так и~внешних, в~том числе не известных заранее. Это означает, 

что реализация ЕИС (ЕИТКС) не должна зависеть от ИС, взаимодействие 

с~которыми она обеспечивает;

     \item[(б)] \textit{использование открытых стандартов доступа}: ЕИС 

(ЕИТКС) должна обеспечивать доступ к~предоставляемым ею сервисам 

посредством единообразных интерфейсов, отвечающих критериям открытых 

стандартов;

     \item[(в)] \textit{технологическая нейтральность}: ЕИС (ЕИТКС) не 

должна требовать от прикладных и~иных ИС использования ка\-кой-ли\-бо 

конкретной технологии, программного обеспечения или аппаратной 

платформы;

     \item[(г)] \textit{стабильность}: неизменность установленных основных 

характеристик, форматов, регламентов функционирования ЕИС (ЕИТКС) 

должна обеспечиваться в~течение необходимого продолжительного времени;

     \item[(д)] \textit{преемственность}: в~случае модернизации ЕИС (ЕИТКС) 

должна обеспечивать продолжительную поддержку функционирования 

необходимых замещаемых протоколов, форматов, спецификаций, 

предоставление сервисов по замещаемым регламентам;

     \item[(е)] \textit{повсеместная доступность}: ЕИС (ЕИТКС) должна 

обеспечивать доступность своих сервисов в~пределах своей территориальной 

деятельности, а~в~особых случаях (по согласованию)~--- за ее пределами;

     \item[(ж)] \textit{постоянная работоспособность}: ЕИС (ЕИТКС) должна 

обеспечивать функционирование и~доступность своих сервисов 

круглосуточно, без перерывов и~отказов в~обслуживании~--- в~пределах 

установленных значений;

     \item[(з)] \textit{нагрузочная способность и~способность 

     к~масштабированию}: ЕИС (ЕИТКС) должна обладать достаточным запасом 

по пропускной способности и~вычислительной нагрузке~--- в~пределах 

установленных значений, в~том числе в~условиях прогнозируемых пиковых 

нагрузок.

     \end{itemize}

     

     Перечисленные требования при выборе архитектуры в~обобщенном 

виде должны рассматриваться как ориентиры при выработке предложений 

и~решений, относящихся к~общесистемным и~системно-техническим вопросам 

построения, создания и~эксплуатации ЕИС (ЕИТКС) различных УЕИП. Эти 

требования должны конкретизироваться, детализироваться и~использоваться 

в процессе технического проектирования ЕИС (ЕИТКС) каждой конкретной 

УЕИП.

     

     \smallskip

     

     6. Определим следующие базовые принципы \textit{разработки 

проектных решений построения и~создания ЕИС (ЕИТКС) конкретной 

УЕИП}:

     \begin{itemize}

     \item \textit{преемственность по отношению к~существующей 

инфраструктуре};

     \item \textit{комплексность автоматизации (информатизации) 

процессов} в~ЕИС (ЕИТКС);

     \item \textit{системность построения} ЕИС (ЕИТКС);

     \item \textit{открытость архитектуры}, во-пер\-вых, обеспечивающая 

взаимодействие сегментов и~компонентов ЕИС (ЕИТКС) друг с~другом через 

стандартные интерфейсы, во-вто\-рых, позволяющая использовать различные 

стандартизованные технические средства;

     \item \textit{обеспечение автоматизации и~оперативности управления} 

ЕИС (ЕИТКС) и~ее составными частями;

     \item \textit{использование математического и~натурного 

моделирования} при создании и~эксплуатации ЕИС (ЕИТКС);

     \item \textit{комплексное обеспечение ИБ} 

     в~ЕИС (ЕИТКС);

     \item \textit{целесообразность и~обоснованность затрат на 

развертывание и~поддержание работоспособности и~актуализации 

автоматизированных ИР} (АИР) в~ЕИС (ЕИТКС);

     \item \textit{направленная эволюция и~максимальное использование 

вложенных инвестиций при дальнейшем развитии и~модернизации} ЕИС 

(ЕИТКС);

     \item \textit{консолидация и~интеграция} как в~узком смысле~--- при 

построении ЕИС (ЕИТКС) и~ее сегментов и~компонентов, так и~в широком 

смысле~--- при выполнении ЕИС (ЕИТКС) функций по обеспечению 

информационного взаимодействия с~различными и~многочисленными УЕИП 

РФ в~электронной форме.

     \end{itemize}

     

     \smallskip

     

     7. В процессе разработки, создания и~развития архитектура ЕИС 

(ЕИТКС) конкретной УЕИП должна \textit{ориентироваться на достижение 

состояния, характеризующегося совокупностью следующих свойств}:

     \begin{itemize}

     \item \textit{адаптивность}: приспособляемость ЕИС (ЕИТКС) 

     к~изменениям внешних условий, в~том числе к~изменениям 

     организационно-функциональной структуры, нормативной базы (или нормативных баз), 

телекоммуникационной среды;

     \item \textit{расширяемость}: обеспечение возможности добавления 

новых или изменения имеющихся функций без изменения остальных 

функциональных компонентов ЕИС (ЕИТКС);

     \item \textit{мобильность (переносимость)}: обеспечение возможности 

переноса специального программного обеспечения и~данных при 

модернизации или замене ап\-па\-рат\-но-программных платформ;

     \item \textit{интероперабельность}: способность к~взаимодействию 

     с~другими информационными системами, не входящими в~состав ЕИС 

(ЕИТКС) (в более широком смысле интероперабельность~--- способность 

программных, информационных или аппаратных ресурсов допускать 

совместное их использование с~другими, заранее не известными при их 

создании ресурсами~\cite{10-zk});

     \item \textit{безопасность}: состояние должной защищенности 

АИР, сегментов и~компонентов 

ЕИС (ЕИТКС).

     \end{itemize}

     

     Приобретение архитектурой ЕИС (ЕИТКС) конкретной УЕИП 

указанных свойств обеспечит выполнение требований, предъявляемых 

к~такому уникальному инфраструктурному образованию.

     

     \smallskip

     

     8.\ В целом ЕИС (ЕИТКС) (включая ее архитектуру) конкретной УЕИП 

должна создаваться как целостная, территориально распределенная (как 

правило), многофункциональная, информационно-коммуникационная 

система, име\-ющая должное документационное обеспечение, в~том числе 

такая ЕИС (ЕИТКС) должна иметь полноценный комплект (комплекты) 

эксплуатационной документации. При этом должны быть две части 

эксплуатационной документации на единую систему: относящейся 

к~внутреннему функционированию и~к функционированию в~составе ЕИП РФ 

как мегасистемы на правах участника (подробнее см.~[1, 2]). 

     

     Как правило, \textit{особенностью создания ЕИС (ЕИТКС) конкретной 

УЕИП является разработка и~ввод ее в~эксплуатацию очередями}. Это 

связано с~существенной сложностью системы (в том числе с~учетом наличия 

унаследованной <<лоскутной>> информатизации), особенностями 

организации и~финансирования работ, с~учетом значительных временн$\acute{\mbox{ы}}$х 

затрат, большой кооперации исполнителей работ.

     

     \smallskip

     

     9. Дадим специально некоторые вводные, более развернутые 

положения, касающиеся системного подхода при определении архитектуры 

ЕИС (ЕИТКС) конкретной УЕИП~--- в~части достижения совокупности 

таких свойств, как адап\-тив\-ность, расширяемость, мобильность 

(переносимость), интероперабельность, безопасность~\cite{10-zk, 11-zk}.

     

     Приобретение ЕИС (ЕИТКС) определенной УЕИП таких свойств, как 

адап\-тив\-ность, расширяемость, мобильность (переносимость), 

интероперабельность, обеспечивается:

     \begin{itemize}

     \item  построением данной системы на принципах открытой 

архитектуры;

     \item использованием при разработке и~эксплуатации средств 

     и~стандартов открытых систем (по возможности). В~том числе следует 

учитывать ряд русифицированных рекомендаций и~стандартов 

международных организаций~\cite{10-zk, 9-zk, 8-zk};

     \item обеспечением взаимодействия сегментов и~их компонентов ЕИС 

(ЕИТКС) различных УЕИП друг с~другом через стандартные интерфейсы, 

поз\-во\-ля\-ющие использовать различные стандартизованные технические 

и~про\-грам\-мные средства. 

     \end{itemize}

     

     Рекомендуется использование в~процессе разработки сегментов 

     и~компонентов ЕИС (ЕИТКС) конкретной УЕИП общего программного 

обеспечения (ОПО) и~технического обеспечения (ТО) ведущих 

производителей (отечественных и~мировых), разработанных на основе 

системы взаимоувязанных международных стандартов с~использованием 

современных технологий и~принципов проектирования. В~определенных 

случаях целесообразно разрабатывать и~внедрять системы 

межведомственных и~ведомственных нор\-ма\-тив\-но-тех\-ни\-че\-ских документов, 

основывающихся на международном опыте в~области информационных 

технологий.

     

     Приобретение ЕИС (ЕИТКС) конкретной УЕИП важнейшего свойства 

<<без\-опас\-ность>>, оцениваемого достижением состояния должной 

защищенности каж\-до\-го сегмента ЕИС в~соответствии с~установленными 

требованиями, обеспечивается учетом при разработке проектных решений 

построения ЕИС (ЕИТКС) \textit{принципов обеспечения информационной 

безопасности}, которые реализуются путем проведения процесса создания 

конкретной системы в~строгом соответствии с~требованиями ГОСТ 

     Р~51583-2000~\cite{11-zk}. С~этой целью в~составе каждого сегмента 

ЕИС (ЕИТКС) создается своя подсистема комплексного обеспечения 

информационной безопасности (ПКОИБ): О-ПКОИБ, К-ПКОИБ, З-ПКОИБ. 

ПКОИБ своего сегмента позволяет обеспечить комплексную защиту 

информации, ИР, персональных данных и~других 

составляющих, передаваемых, накапливаемых и~обрабатываемых в~каждом 

сегменте ЕИС (ЕИТКС) путем адекватной защиты от угроз. В~свою очередь, 

адекватность защиты от угроз для каждого сегмента ЕИС (ЕИТКС) должна 

быть определена в~согласованном с~ФСБ России и~Гостехкомиссией 

(\mbox{ФСТЭК}) России и~утвержденном документе <<Модель нарушителя 

и~модель угроз>>. Также должны выполняться установленные требования по 

защите информации и~персональных данных для каждого сегмента ЕИС 

(ЕИТКС). Работы по созданию О-ПКОИБ, К-ПКОИБ, З-ПКОИБ должны 

проводиться в~соответствии с~федеральными законами, стандартами, 

действующими руководящими и~нормативными документами 

уполномоченных органов исполнительной власти, включая нормативные 

документы ФСБ России и~ФСТЭК России.

     

     Сегменты ЕИС (ЕИТКС), обрабатывающие информацию 

     и~персональные данные, должны быть защищены и~аттестованы по 

требованиям безопасности информации в~соответствии с~действующими 

руководящими документами ФСБ России и~ФСТЭК России. Объектами 

защиты в~сегментах ЕИС (ЕИТКС) являются программно-технические 

комплексы (ПТК) различного назначения, размещаемые на объектах ЕИС 

(ЕИТКС) определенной УЕИП.

     

     \smallskip

     

     10. \textit{Принципы системности построения и~открытости 

архитектуры} пред\-усмат\-ри\-вают, в~частности, возможность развития ЕИС 

(ЕИТКС) конкретных УЕИП в~следующих направлениях~\cite{6-zk}:

     \begin{itemize}

     \item \textit{увеличение количества пользователей} ЕИС (ЕИТКС);

     \item \textit{расширение состава внешних ИС} (АИС), 

взаимодействующих с~ЕИС (ЕИТКС) данной УЕИП;

     \item \textit{расширение возможностей контроля и~управления 

функционированием сис\-те\-мы};

     \item \textit{повышение уровня унификации и~стандартизации 

     в~сис\-теме};

     \item \textit{повышение производительности сис\-те\-мы};

     \item \textit{совершенствование форм представления информации, 

используемых в~сис\-теме};

     \item \textit{консолидация, интеграция и~централизация 

ИР в~сис\-теме}.

     \end{itemize}

     

     \smallskip

     

     11. Создание и~эффективное функционирование ЕИС (ЕИТКС) 

конкретной УЕИП~--- это сложные процессы, если принять во внимание 

сложность системы: ее, как правило, территориальную распределенность 

и~взаимодействие с~большим числом разнородных ИС (АИС) других органов 

и~организаций. Поэтому необходимо использовать \textit{методы 

математического и~натурного моделирования}.

     

     Математическое моделирование целесообразно использовать для 

определения характеристик системы, оценки и~прогнозирования поведения 

сегментов и~их компонентов в~различных ситуациях, совершенствования 

архитектуры, определения характеристик настройки и~требований 

к~сегментам и~компонентам системы, оценки уровня качества и~определения 

путей повышения эф\-фек\-тив\-ности функционирования и~эксплуатации. Может 

быть также разработан комплект математических моделей разного 

назначения, например: математическая модель телекоммуникационной 

системы ЕИС (ЕИТКС), модель надежности системы, сервисная модель 

эксплуатации и~сопровождения системы и~др.

     

     Следует в~обязательном порядке отработать для ЕИС (ЕИТКС) 

\textit{модель создания профиля среды} как открытой системы в~соответствии 

с~принятой стратегией и~согласно рекомендациям руководства 

     Р~50.1.041-2002~\cite{8-zk}. При этом применительно к~ЕИС (ЕИТКС) 

должны быть выполнены последовательно пять этапов моделирования 

системы: область действия, анализ требований, логический проект, 

физический проект, эксплуатационный проект (подробнее см.\ в~\cite{2-zk}).

     

     В ЕИС (ЕИТКС) необходимо создать постоянно действующий 

ими\-та\-ци\-он\-но-от\-ла\-до\-чный комплекс и~разработать соответствующие 

имитаторы с~целями:

     \begin{itemize}

     \item проведения натурного моделирования;

     \item отладки и~испытаний сегментов ЕИС (ЕИТКС) соответствующей 

УЕИП и~их компонентов;

     \item проверки соответствия требованиям нормативных документов 

     и~сертификации сегментов ЕИС (ЕИТКС) и~их компонентов (в~том числе для 

имитации отсутствующих компонентов системы);

     \item предварительной отработки различных изменений, планируемых 

к внесению их в~конфигурацию действующей ЕИС (ЕИТКС) 

соответствующей УЕИП, и~др.

     \end{itemize}

     

     \smallskip

     

     12. \textit{Учет принципа обеспечения автоматизации управления 

инфраструктурой} ЕИС (ЕИТКС) соответствующей УЕИП при разработке 

проектных решений предусматривает \textit{выполнение следующих 

функций}:

\begin{itemize}

     \item \textit{организации автоматизированного сбора и~хранения 

данных об аппаратной и~программной конфигурации сегментов} ЕИС 

(ЕИТКС) и~их компонентов;

     \item \textit{автоматизированного управления} ЕИС (ЕИТКС) данной 

УЕИП в~целом и~ее сегментами;

     \item \textit{диагностики состояния и~поиска неисправностей}, 

включая обеспечение инструментальными средствами, для эффективного 

анализа аппаратных и~программных проблем, возникающих на удаленных 

ПТК из состава единой системы соответствующей УЕИП;

     \item \textit{обеспечения возможности удаленного администрирования 

объектовых ПТК} из состава ЕИС (ЕИТКС), включая управление и~контроль 

установки, обновления и~выполнения программ;

     \item \textit{создание отчетов о работоспособности 

     и~производительности вычислительных ресурсов} ЕИС (ЕИТКС).

     \end{itemize}

     

     \smallskip

     

     13. \textit{Соблюдение принципа обеспечения 

ИБ} при построении ЕИС (ЕИТКС) соответствующей УЕИП 

следует реализовать:

     \begin{enumerate}[(1)]

     \item за счет \textit{привлечения специализированной 

организации в~части обеспечения ИБ}. Эта 

организация должна участвовать в~процессе создания ЕИС (ЕИТКС) на всех 

ключевых этапах в~соответствии с~требованиями ГОСТ 

     Р~51583-2000~\cite{11-zk}. У~специализированной организации 

должны быть все необходимые лицензии ФСБ России и~ФСТЭК России для 

проведения работ по данному направлению, а также опыт работ по 

обеспечению ИБ при создании крупных территориально распределенных 

ИС;

     \item за счет \textit{понимания важности данной работы 

по ИБ} и~всесторонней помощи на уровне заказчика ЕИС (ЕИТКС).

     \end{enumerate}

     

     Существенным условием гарантированности соблюдения принципа 

обеспечения ИБ является и~тот факт, что с~учетом стратегического значения 

для России создания ЕИС (ЕИТКС) весь процесс построения данной системы 

находится под постоянным контролем уполномоченных организаций в~части 

безопасности и~технической защиты информации~--- ФСБ России и~ФСТЭК 

России.

     

     \smallskip

     

     14. \textit{Выполнение принципа экономической целесообразности 

     и~обосно\-ван\-ности финансовых затрат на развертывание и~поддержание 

работоспособ\-ности и~актуализации} ЕИС (ЕИТКС) оценивается 

сокращением затрат (единая система вместо многочисленных 

самостоятельных ИС), а также повышением качества информационной 

поддержки в~сфере деятельности конкретной УЕИП. С другой стороны, 

указанный принцип обеспечивается сокращением затрат на развертывание 

и~поддержание работоспособности и~актуализации компонентов ЕИС (ЕИТКС) 

этой УЕИП путем типизации процессов, унификации решений 

и~оптимизации состава данной системы.

     

     \smallskip

     

     15. \textit{Соблюдение принципов преемственности, направленной 

эволюции и~максимального использования вложенных инвестиций} следует 

считать обязательным при создании широкомасштабных, функционально 

сложных компонентов ЕИС (ЕИТКС) соответствующей УЕИП с~длительным 

жизненным циклом, с~обеспечением поэтапного создания очередями 

и~дальнейшего развития и~модернизации как компонентов и~в целом данной 

системы, так и~ИС (АИС), взаимодействующих с~ней~\cite{6-zk}.

     

     \smallskip

     

     16. Заключая рассмотрение требований и~предложений по системному 

подходу к~ЕИС конкретной УЕИП и~ее 

архитектуре, необходимо:

     \begin{enumerate}[(1)]

          \item акцентировать внимание на \textit{особой значимости 

решения организационных вопросов} при проектировании, создании 

и~эксплуатации ЕИС (ЕИТКС) в~течение ее жизненного цикла;

     \item так же \textit{четко, пунктуально отслеживать 

соответствующие решения, планы и~графики их выполнения}.

     \end{enumerate}

     

     Важное положение организационного плана: ЕИС (ЕИТКС) 

соответству\-ющей УЕИП должна быть наделена правами координации, 

формирования и~развития ИР в~сфере своей 

деятельности и~включения в~ЕИП РФ

своих ИР.

     

\section{Заключение}



\vspace*{-2pt}



     Статья расширяет и~конкретизирует системные подходы 

к~формированию нор\-ма\-тив\-но-технической базы в~части требований 

к~архитектурному построению ЕИС

организаций~--- участников ЕИП РФ с~учетом всех этапов жизненного цикла.

     

     Предложены обобщенные требования к~системам и~базовые принципы 

разработки проектных решений по их созданию. 

     

     Определены необходимые свойства систем, направления их 

возможного развития, состав необходимого документационного обеспечения.

     

     Показана необходимость соблюдения принципов преемственности, 

эволюционности, обеспечения ИБ при 

построении ЕИС, экономической целесообразности и~обоснованности 

финансовых затрат на развертывание и~поддержание работоспособности.

     

     Сформулированы основные положения по использованию методов 

математического и~натурного моделирования на всех этапах жизненного 

цикла ЕИС.

     

     Подчеркнута особая значимость решения организационных вопросов 

при проектировании, создании и~эксплуатации ЕИС (ЕИТКС) различных 

УЕИП на всех этапах жизненного цикла.

     

     Предложенные в~статье подходы могут стать методической основой 

при формировании нормативно-технической базы создания 

ЕИС различных УЕИП в~части их архитектурного 

построения.



\vspace*{-6pt}

     

{\small\frenchspacing

{%\baselineskip=10.8pt

\addcontentsline{toc}{section}{Литература}





\begin{thebibliography}{99}



\vspace*{-2pt}



      \bibitem{1-zk}

      \Au{Зацаринный А.\,А., Киселев Э.\,В.} Некоторые подходы к~формированию 

нормативно-тех\-ни\-че\-ской базы для создания единого информационного пространства 

России~/\!\!/ Системы и~средства информатики, 2014. Т.~24. №\, 4. С.~206--220.

      \bibitem{2-zk}

      \Au{Зацаринный А.\,А., Киселев Э.\,В.} Некоторые подходы к~формированию 

нормативно-тех\-ни\-че\-ской базы единого информационного пространства России в~части 

информационных ресурсов~/\!\!/ Системы и~средства информатики, 2015. Т.~25. №\,1. 

С.~157--169.

      \bibitem{3-zk}

      Электронная Россия: Федеральная целевая программа. Последние изменения 

2008~г.

      \bibitem{4-zk}

      \Au{Штрик А.\,А.} Использование информационно-коммуникационных 

технологий для экономического развития и~государственного управления в~странах 

современного мира~/\!\!/ Информационные технологии. Приложение, 2009. №\,6. 32~с.

      \bibitem{5-zk}

      \Au{Штрик А.\,А.} Электронные технологии в~деятельности органов 

государственной власти России: анализ и~перспективы развития~/\!\!/ Информационные 

технологии. Приложение, 2009. №\,10. 32~с.

      \bibitem{6-zk}

      \Au{Зацаринный А.\,А., Ионенков Ю.\,С., Козлов~С.\,В.} Некоторые вопросы 

проектирования информационно-телекоммуникационных систем.~--- М.: 

ИПИ РАН, 2010. 218~с.

      \bibitem{7-zk}

      \Au{Зацаринный А.\,А., Шабанов А.\,П.} Технология информационной поддержки 

деятельности организационных систем на основе ситуационных центров.~--- М.: ТОРУС 

ПРЕСС, 2015. 235~с.

      \bibitem{10-zk} %8

      Стандарты на открытые системы серии ГОСТ Р~ISO/IEC ТО 10000-99. 

Информационная технология. Основы и~таксономия международных функциональных 

стандартов ВОС.

      \bibitem{11-zk} %9

      ГОСТ Р 51583-2000. Защита информации. Порядок создания автоматизированных 

систем в~защищенном исполнении. 



      \bibitem{9-zk} %10

      Р 50.1.022-2000. Государственный профиль взаимосвязи открытых систем России. 

Информационная технология.



      \bibitem{8-zk} %11

      Информационные технологии. Руководство по проектированию профилей среды 

открытой системы (СОС) организации-пользователя: Рекомендации по стандартизации 

РФ Р 50.1.041-2002. Введены в~действие с~1~января 2004~г.



   \end{thebibliography}

} }





\vspace*{-4pt}



\hfill{\small\textit{Поступила в~редакцию 16.07.15}}









\vspace*{7pt}



\hrule



\vspace*{2pt}



\hrule



%\vspace*{9pt}





\def\tit{REGARDING SOME APPROACHES TO CREATION OF~THE~REGULATORY 

AND~TECHNICAL BASE CONCERNING REQUIREMENTS TO~ARCHITECTURAL 

DEVELOPMENT OF~INFORMATION SYSTEMS OF~RUSSIAN INFORMATION SPACE 

MEMBERS}





\def\titkol{Regarding some approaches to creation of~the~regulatory 

and~technical base} % concerning requirements to~architectural  development of~information systems of~Russian information space  members}





\def\aut{A.\,A.~Zatsarinny and E.\,V.~Kiselev}

\def\autkol{A.\,A.~Zatsarinny and E.\,V.~Kiselev}







\titel{\tit}{\aut}{\autkol}{\titkol}



\vspace*{-6pt}



\noindent

Institute of Informatics Problems, 

Federal Research Center ``Computer Science and

Control'' of the Russian Academy of Sciences, 44-2 Vavilov Str.,

Moscow 119333, Russian Federation







\def\leftfootline{\small{\textbf{\thepage}

\hfill Sistemy i Sredstva Informatiki~---

Systems and Means of Informatics 2015 vol~25 no\,3}

}%

 \def\rightfootline{\small{Sistemy i Sredstva Informatiki~---

 Systems and Means of Informatics   2015   vol~25   no\,3

\hfill \textbf{\thepage}}}   

      

  

  \vspace*{2pt} 

  



\Abste{The article continues the topic about regulatory and technical base forming for creation 

of information systems of organizations which are members of the Russian Federation Unified 

Information Space (RF UIS). The article shows generalized requirements to systems including 

main principles of systems creation project development, systems necessary properties, 

necessary documentation, potential development of information systems, requirements to 

mathematics and natural modeling during systems creation and operation, requirements to 

automated management of systems working, providing of information security, rationality in 

finances of project during system roll-out,  maintenance and upgrade, providing of continuity, 

directed evolution  and efficient use of investments, emphasized value of organization problem 

solution during the lifecycle. Propositions of the article are oriented to practical using as a 

science and methodical base for creation of regulatory and science base of RF UIS members 

information systems development.}           

 

\KWE{Russian Federation Unified Information Space (RF UIS); RF UIS member organization; 

systemic approach; UIS member system architecture; secured information resources; 

centralization; integration; information security; telecommunication provision}   



\DOI{10.14357\!/\!08696527150311}



\vspace*{-4pt}



%\Ack







\renewcommand{\bibname}{\protect\rmfamily References}

%\renewcommand{\bibname}{\large\protect\rm References}



{\small\frenchspacing

{%\baselineskip=10.8pt

\addcontentsline{toc}{section}{References}

\begin{thebibliography}{99}

      

\bibitem{1-zk-1}

\Aue{Zatsarinny, A.\,A., and E.\,V.~Kiselev}. 2014. Nekotorye podkhody k~formirovaniyu 

normativno-tekhnicheskoy bazy dlya sozdaniya edinogo informatsionnogo prostranstva Rossii 

[Regarding some approaches to creation of regulatory and technical base for Russian Unified 

Informational Space creation]. \textit{Sistemy i~Sredstva Informatiki}~--- \textit{Systems and 

Means of Informatics}  24(4):206--220.

\bibitem{2-zk-1}

\Aue{Zatsarinny, A.\,A., and E.\,V.~Kiselev}. 2015. Nekotorye podkhody k formirovaniyu 

normativno-tehnicheskoy bazy edinogo informatsionnogo prostranstva Rossii v~chasti 

informatsionnykh resursov [Regarding some approaches to creation of regulatory and technical 

base of Russian Unified Informational Space in the field of information resources]. 

\textit{Sistemy i~Sredstva Informatiki}~--- \textit{Systems and Means of Informatics}  

25(1):157--169.

\bibitem{3-zk-1}

Elektronnaya Rossiya: Federal`naya tselevaya programma. Poslednie izmeneniya 2008~g. 

[Electronic Russia. Federal objective program. Update of~2008].

\bibitem{4-zk-1}

\Aue{Shtrik, A.\,A.} 2009. Ispol'zovanie informatsionno-kommunikatsionnykh tekhnologiy dlya 

ekonomicheskogo razvitiya i~gosudarstvennogo upravleniya v~stranakh sovremennogo mira 

[Information and telecommunication technology application to economic development and 

public administration of modern world countries]. \textit{Informatsionnye Tekhnologii. 

Prilozhenie} [Appendix to Information Technologies]. Vol.~6. 32~p.

\bibitem{5-zk-1}

\Aue{Shtrik, A.\,A.} 2009. Elektronnye tekhnologii v~deyatel`nosti organov gosudarstvennoy 

vlasti Rossii: Analiz i~perspektivy razvitiya  [Electronic technologies in Russian authorities 

activities]. \textit{Informatsionnye Tekhnologii. Prilozhenie} [Appendix to Information 

Technologies]. Vol.~10. 32~p.

\bibitem{6-zk-1}

\Aue{Zatsarinny, A.\,A., Y.\,S.~Ionenkov, and S.\,V.~Kozlov}. 2010. \textit{Nekotorye 

voprosy proektirovaniya informatsionno-telekommunikatsionnykh system} [Some issues of 

information and telecommunication system designing]. Moscow: IPI RAN. 218~p.

\bibitem{7-zk-1}

\Aue{Zatsarinny, A.\,A., and A.\,P.~Shabanov}. 2015. \textit{Tekhnologiya informatsionnoy 

podderzhki deyatel'nosti organizatsionnykh sistem na osnove situatsionnykh centrov} 

[A~technology of information support of organization systems activities based on situation 

centers]. Moscow: TORUS PRESS. 235~p.



\bibitem{10-zk-1} %8

GOST R ISO/IEC TO 10000-99. Standarty na otkrytye sistemy. Informatsionnaya tekhnologiya. 

Osnovy i~taksonomiya mezhdunarodnykh funktsional`nykh standartov\linebreak

%\pagebreak

%\noindent

 VOS [Open system's 

standards. Information technology. Foundations and taxonomy of

 VOS international functional  standards].

\bibitem{11-zk-1} %9

GOST R 51583-2000. Zashchita informatsii. Poryadok sozdaniya avtomatizirovannykh sistem 

v~zashchishchennom ispolnenii [Information security. An order of secured automated 

information systems creation]. 



\bibitem{9-zk-1} %10

R 50.1.022-2000. Gosudarstvenny profil vzaimosvyazi otkrytykh sistem Rossii. 

Informatsionnaya tekhnologiya [A~state profile of interconnection of Russian open systems. 

Information technology].





\bibitem{8-zk-1} %11

RF R 50.1.041-2002. Actual from 01.01.2004. Rekomendatsii po standartizatsii. 

Informatsionnye tekhnologii. Rukovodstvo po proektirovaniyu profiley sredy 

otkrytoy sistemy 

(SOS) organizatsii-pol'zovatelya [Standartization recommendations. Information technologies. 

Guidance on design of user-organization open system environment profiles]. 







   \end{thebibliography}

} }





\vspace*{-1pt}



\hfill{\small\textit{Received July 16, 2015}}   

      

      

      \Contr

      

      \noindent

      \textbf{Zatsarinny Alexander A.} (b.\ 1951)~--- Doctor of Science in technology, 

professor, Deputy Director, Federal Research Center 

``Computer Science and Control'' of the Russian Academy of Sciences, 44-2 Vavilov Str., Moscow 

119333, Russian Federation; azatsarinny@ipiran.ru

      

      \vspace*{3pt}

      

\noindent

\textbf{Kiselev Eduard V.} (b.\ 1938)~---

senior scientist, Institute of Informatics Problems, Federal Research Center ``Computer Science and 

Control'' of the Russian Academy of Sciences, 44-2 Vavilov Str., Moscow 119333, Russian Federation;\linebreak

 ekiselev@ipiran.ru

      

 \label{end\stat}





\renewcommand{\bibname}{\protect\rm Литература} 

